% 图7-2：检测信号产生候选解释，独立证据才形成裁决。
\begin{tikzpicture}[
  x=1mm,
  y=1mm,
  font=\figurefont\footnotesize,
  text=BookInk,
  signal/.style={draw=BookAmber,fill=BookAmber!9,line width=1.1pt,rounded corners=1.2mm,text width=24mm,minimum height=19mm,align=center,inner sep=1.8mm},
  explanation/.style={draw=BookMuted,fill=BookPaper,line width=0.8pt,rounded corners=1.2mm,text width=30mm,minimum height=12mm,align=center,inner sep=1.5mm},
  check/.style={draw=FigureInput,fill=FigureInput!7,line width=0.8pt,text width=30mm,minimum height=12mm,align=center,inner sep=1.5mm},
  verdict/.style={draw=FigureOutput,fill=FigureOutput!6,line width=0.8pt,text width=23mm,minimum height=12mm,align=center,inner sep=1.4mm},
  candidate/.style={-{Stealth[length=1.9mm]},draw=BookAmber,line width=0.8pt,dashed,rounded corners=1mm},
  evidence/.style={-{Stealth[length=1.9mm]},draw=BookBlue,line width=0.9pt},
  decision/.style={-{Stealth[length=1.9mm]},draw=BookTeal,line width=1.0pt},
  edge label/.style={fill=white,inner sep=0.8pt,font=\figurefont\scriptsize,text=BookMuted,align=center}
]
  \node[signal] (signal) at (15,20) {高损失或\\预测分歧\\{\scriptsize 检测线索}};

  \node[explanation] (wrong) at (58,40) {候选解释一\\知识确有错误};
  \node[explanation] (rare) at (58,20) {候选解释二\\样本正确但罕见};
  \node[explanation] (model) at (58,0) {候选解释三\\模型能力或分布不匹配};

  \node[check] (wrong-check) at (108,40) {核对原始对象、规范\\与独立参考证据};
  \node[check] (rare-check) at (108,20) {检查类别频率、覆盖\\与专业复核结果};
  \node[check] (model-check) at (108,0) {更换模型或输入条件\\比较留出行为};

  \node[verdict,draw=BookAmber,fill=BookAmber!9] (repair) at (151,40) {确认错误\\进入修正};
  \node[verdict] (keep) at (151,20) {知识可保留\\维护稀有覆盖};
  \node[verdict,draw=BookMuted,fill=white] (limit) at (151,0) {不能据此改标\\限制检测解释};

  \draw[candidate] (signal.east) -- (wrong.west);
  \draw[candidate] (signal.east) -- (rare.west);
  \draw[candidate] (signal.east) -- (model.west);
  \node[edge label,anchor=south] at (34,22) {待核查解释};

  \draw[evidence] (wrong) -- (wrong-check);
  \draw[evidence] (rare) -- (rare-check);
  \draw[evidence] (model) -- (model-check);
  \draw[decision] (wrong-check) -- (repair);
  \draw[decision] (rare-check) -- (keep);
  \draw[decision] (model-check) -- (limit);

  \node[anchor=north,font=\figurefont\scriptsize,text=BookMuted,align=center] at (84.5,-11) {同一个排序分数不能区分三种原因；裁决来自右侧检查，而不是来自分数本身。};
\end{tikzpicture}
